% Generated draft: requires mathematical and source review.
\documentclass[11pt]{article}
\usepackage[margin=25mm]{geometry}
\usepackage{fontspec}
\setmainfont{Latin Modern Roman}
\usepackage{amsmath,amssymb,mathtools,unicode-math}
\setmathfont{Latin Modern Math}
\usepackage{enumitem,needspace,xurl,hyperref}
\setlength{\parindent}{0pt}
\setlength{\parskip}{5pt}
\setlength{\emergencystretch}{3em}
\newcommand{\N}{\mathbb{N}}
\newcommand{\Z}{\mathbb{Z}}
\newcommand{\Q}{\mathbb{Q}}
\newcommand{\R}{\mathbb{R}}
\newcommand{\C}{\mathbb{C}}
\newcommand{\F}{\mathbb{F}}
\newcommand{\A}{\mathbb{A}}
\newcommand{\PP}{\mathbb{P}}
\newcommand{\E}{\mathbb{E}}
\newcommand{\eps}{\varepsilon}
\newcommand{\dd}{\,\mathrm{d}}
\begin{document}
\section*{UnsolvedMath ID: 800017}
\textbf{Source problem code:} catalogue entry
\subsection{Definitions and mathematical statement}
A semigroup is a nonempty set \(S\) equipped with an associative binary operation. It is finitely generated if there is a finite nonempty subset \(A\subseteq S\) such that every \(s\in S\) can be written as \(s=a_1\cdots a_n\) for some integer \(n\geq 1\) and elements \(a_1,\ldots,a_n\in A\). An element \(e\in S\) is idempotent if \(e^2=e\); thus having no idempotents means
\[
\forall e\in S,\qquad e^2\neq e.
\]
\textbf{Problem (A. A. Borisov).} Describe the finitely generated Lesokhin semigroups having no idempotents. Interpreting ``describe'' as a request for classification up to semigroup isomorphism, the target is the class
\[
\mathcal C=\{S\mid S\text{ is a finitely generated Lesokhin semigroup and }(\forall e\in S)\ e^2\neq e\}.
\]
\textbf{Unresolved qualification.} The supplied record does not define ``Lesokhin semigroup.'' That term remains an unspecified source-dependent condition here, so the displayed class is not a self-contained mathematical specification. The source also does not specify the desired form of the description.

\textbf{Supplied background assertion.} The record attributes to reference [Borisov4] the assertion that every finitely generated \emph{separative} Lesokhin semigroup without idempotents is isomorphic to a subsemigroup of an infinite monogenic semigroup. The latter may be represented as
\[
M=\{a^n:n\in\mathbb Z_{\geq 1}\},\qquad a^m a^n=a^{m+n},
\]
with all powers distinct; equivalently, \(M\cong(\mathbb Z_{\geq 1},+)\). Thus the supplied background asserts, under its additional separativity hypothesis, the existence of an injective map \(\varphi:S\to\mathbb Z_{\geq 1}\) satisfying
\[
\varphi(xy)=\varphi(x)+\varphi(y)\qquad(x,y\in S).
\]
This is reported as an existing result in the source, not as the requested classification or as an independently verified theorem. Neither the definition of ``separative'' nor the bibliographic entry [Borisov4] is included in the supplied record; the exact hypotheses and wording of that assertion require review.
\subsection{Short English statement}
The problem, attributed to A. A. Borisov, asks for a description of all finitely generated Lesokhin semigroups with no idempotents. The supplied background reports an embedding result for the additional separative case. A self-contained formulation requires the missing definition of a Lesokhin semigroup and clarification of the intended classification.
\subsection{Sources}
\begin{itemize}
\item Local UnsolvedMath registry: permanent ID \(800017\), source problem code \texttt{catalogue entry}, title \emph{The Lyapin's Notebook --- Problem 17}. The supplied registry labels the problem \emph{open}; this is supplied curation, not an independently verified status assessment.
\item Bershadsky, Kublanovsky, Mashevitzky, \emph{The Lyapin's Notebook: Unsolved Problems in Semigroup Theory} (2026), Problem 17. Supplied source URL: \url{https://arxiv.org/abs/2604.04763}. The supplied provenance records access on 2026-10-16 and extraction by \texttt{arxiv\_latex}; these are source-record provenance claims, not a claim of access or verification for this draft.
\item The problem text attributes the question to A. A. Borisov. The registry's proposer field lists Bershadsky, Kublanovsky, Mashevitzky and its proposed-year field gives 2026; the original proposal date is not established by the supplied evidence.
\item Reference [Borisov4] is cited in the supplied text for the separative-case assertion. Its author details, title, publication information, and URL are not supplied, so the reference and its exact result cannot be checked from this record alone.
\end{itemize}
\end{document}
